\chapter{\texttt{tricycle\_env\_cfg.py}: the environment configuration}
\label{ch:env_cfg}

\textbf{Role:} pure data. It describes the scene (ground, light, car, hidden
slab, soil), the physics (Newton with a coupled rigid + MPM solver), the
spaces, the time steps and the termination thresholds. It contains no
behaviour except two methods on the cfg class: \code{__post_init__}, which
assembles the physics config, and \code{play_mode}, which adjusts things for
playback. 300 lines. It is the file where almost every hard-won fact from the
runbook lives, so it gets the longest chapter.

\begin{isaacbox}{\code{configclass}}
\code{@configclass} (\srcpath{source/isaaclab/isaaclab/utils/configclass.py}) is
Isaac Lab's wrapper around Python's \code{dataclass}. It lets you write
fields without type annotations, makes mutable defaults safe by deep-copying
them per instance, adds \code{to_dict}, \code{from_dict}, \code{replace}, \code{copy} and \code{validate},
and calls \code{__post_init__} after the instance is built. A field set to
\code{MISSING} must be filled in by a subclass or \code{validate()} raises.
Everything in this file is a configclass instance; nothing is executed by
the physics engine until \srcpath{DirectRLEnv.__init__} reads it.
\end{isaacbox}

\section{Header}

\codepart{\srcroot/luna_hifi_tasks/tricycle/tricycle_env_cfg.py}{1}{21}

\begin{explain}
\lis{1}{4} What the task is.
\lis{6}{16} The design is copied from Isaac Lab's own coupled task,
  \srcpath{source/isaaclab_tasks/isaaclab_tasks/contrib/ur10_particle_push/ur10_particle_push_env_cfg.py},
  which pushes particles with a robot paddle. The six bullets are the rules
  that task obeys and that this file obeys too; each reappears below next to
  the line that implements it.
\lis{18}{21} Where the car comes from and why the MJCF cannot be loaded
  directly at run time (Chapter~\ref{ch:asset}).
\end{explain}

\section{Imports}

\codepart{\srcroot/luna_hifi_tasks/tricycle/tricycle_env_cfg.py}{23}{49}

\begin{explain}
\li{23} Postponed evaluation of annotations, so type hints can name classes
  freely.
\li{25} \code{pathlib.Path} for locating the asset relative to this file
  (line 83).
\li{27} \code{isaaclab.sim} is imported as \code{sim_utils} by Isaac Lab
  convention; it holds the spawner configs (\code{UsdFileCfg},
  \code{CuboidCfg}, \code{GroundPlaneCfg}, \code{DomeLightCfg}).
\li{28} \code{ImplicitActuatorCfg}: a PD drive executed inside the physics
  solver (Chapter~\ref{ch:physics}).
\li{29} The three asset config types used in the scene: an articulated robot,
  a plain rigid body, and a static scene element with no physics state of its
  own.
\li{30} The base class for direct-workflow environment configs.
\li{31} The base class for scene configs.
\li{32} The simulation config (time step, physics backend, visualisers).
\li{33} \code{UsdPhysicsRigidBodyCfg}: the schema config that marks a prim as
  a rigid body and can make it kinematic.
\li{34} \code{RigidBodyMaterialBaseCfg}: friction values for a rigid body.
\li{35} The \code{configclass} decorator. The runbook notes that this exact
  import path matters: an older \code{from isaaclab.utils import configclass}
  form, combined with stale \code{__pycache__}, once produced a confusing
  ``\code{'module' object is not callable}''.
\li{37} \code{MPMObjectCfg}: the asset config for a particle body.
\lis{38}{43} The Newton physics configs: the MuJoCo-Warp rigid solver, the
  MPM solver, the top-level \code{NewtonCfg}, and the collision-pipeline
  config.
\li{44} \code{NewtonCollisionPropertiesCfg}: per-prim Newton collision
  settings (margin, gap), used on the hidden slab.
\li{45} The particle spawner (\code{MPMGridCfg}) and its material
  (\code{MPMParticleMaterialCfg}).
\li{47} The coupler configs from the \code{isaaclab_contrib} package: an
  entry, the proxy coupler, and a proxy mapping
  (Section~\ref{sec:coupling}).
\li{49} From our own model file: the chassis height and the joint names.
  Nothing else from \code{tricycle.py} is needed at run time.
\end{explain}

\section{Constants}

\codepart{\srcroot/luna_hifi_tasks/tricycle/tricycle_env_cfg.py}{50}{74}

\begin{explain}
\lis{54}{55} The names of the two coupler entries. They are only used to
  link the proxy mapping to its source and destination (lines 272--273), but
  naming them once avoids typos.
\li{57} The MPM grid spacing, 5~cm. It sets soil resolution, the number of
  particles (line 62--63) and the collider margin below.
\li{58} Half a voxel. Used both as the collider margin on the slab (line 152)
  and as the vertical offset of the soil box (lines 67--68), so the lowest
  particle layer starts half a cell above the slab surface, where the
  implicit solver can see the collider before the particles touch it.
\li{59} One particle per grid cell along each axis. Higher values multiply
  the particle count by the cube of the factor.
\li{60} Soil colour (RGB, 0--1), used for the spawner and for the particle
  colour in the viewer.
\lis{62}{65} The strip: 1.5~m long, 0.6~m wide, 6~cm deep, starting 40~cm
  behind the car so that the car spawns on it with room behind. With the
  resolution rule from Chapter~\ref{ch:physics},
  $\lceil 1.5/0.05\rceil \times \lceil 0.6/0.05\rceil \times \lceil 0.06/0.05\rceil = 30\times12\times2 = 720$
  particles per env.
\li{67} The lower corner of the particle box: rear end of the strip, left
  edge, half a voxel above $z=0$.
\li{68} The upper corner: front end, right edge, strip depth above the lower
  $z$. Because \code{MPMGridCfg} coordinates are local to the asset prim and
  the asset's \code{init_state} is left at the origin (line 120), these are
  also env-local coordinates.
\lis{70}{72} The rationale for the slab (also a green box in
  Chapter~\ref{ch:physics}).
\li{73} Slab size: the strip plus a 1~m margin in $x$ and $y$, 10~cm thick.
\li{74} Slab centre: under the middle of the strip, 5~cm below $z=0$, so its
  top face lies exactly at $z=0$.
\end{explain}

\section{The car}

\codepart{\srcroot/luna_hifi_tasks/tricycle/tricycle_env_cfg.py}{76}{112}

\begin{explain}
\lis{80}{82} Comment: where the asset is and when to regenerate it.
\li{83} The repository root, found by walking up from this file:
  \code{parents[0]} is \srcpath{tricycle/}, \code{parents[1]} is
  \srcpath{luna_hifi_tasks/}, \code{parents[2]} is the repository. This
  replaced a hard-coded \texttt{C:\textbackslash ...} path, so the repository can be cloned
  anywhere.
\li{84} The USD entry file.
\li{86} An \code{ArticulationCfg} describes a robot: where to put it, what to
  spawn, its initial state, and its actuators.
\li{87} The prim path uses the \code{{ENV_REGEX_NS}} macro, which
  \code{InteractiveScene} expands to \srcpath{/World/envs/env_.*}. One car is
  spawned under each env prim (Section~\ref{sec:prim_paths}).
\li{88} \code{UsdFileCfg} spawns a USD file by reference. \code{usd_path}
  must be a string, hence \code{str(...)}.
\li{89} Initial (and reset) state.
\li{90} Spawn position relative to the env origin: $x=y=0$, and $z$ equal to
  the chassis height that puts the wheel bottoms at $z=0$, plus the soil
  offset, plus the soil depth, plus 1~cm clearance. The car therefore starts
  just above the soil surface and settles onto it in the first few steps.
\li{91} Orientation as a quaternion in Isaac Lab's $(w, x, y, z)$ order;
  identity.
\lis{92}{93} All joints start at zero position and zero velocity. The keys
  are regexes over joint names; \code{".*"} matches all.
\li{95} The actuator groups. Each key is a label; each value binds one
  actuator model to a set of joints selected by name.
\li{96} The steering group.
\li{97} Bind it to the joint named \code{steer}.
\li{98} Maximum torque the solver may apply, 5~N\,m.
\li{99} Maximum joint speed, 5~rad/s.
\li{100} PD stiffness 20~N\,m/rad: the drive pulls the fork towards the
  commanded angle.
\li{101} PD damping 1~N\,m\,s/rad.
\li{102} End of the steering group.
\li{103} The rear-wheel group, bound to both \code{rear_left} and
  \code{rear_right} (line 104).
\li{105} Torque limit 2~N\,m per wheel. With a 10~cm radius that is 20~N of
  tractive force per wheel, about 60\% of the 6.6~kg car's weight in total, a
  plausible traction limit for soil.
\li{106} Speed limit 30~rad/s, above the 20~rad/s the policy can ask for.
\li{107} Stiffness zero: there is no position target, so the drive is a pure
  \emph{velocity} controller.
\li{108} Damping 0.5~N\,m\,s/rad: the torque is $0.5\times(\text{target speed}
  - \text{actual speed})$, saturating at the 2~N\,m limit.
\li{110} Comment: the front wheel has no actuator; it rolls freely with the
  MJCF's small joint damping.
\lis{111}{112} Close the dictionary and the cfg.
\end{explain}

\begin{isaacbox}{how a PD implicit actuator becomes a velocity drive}
An implicit drive computes torque $=$
\code{stiffness}$\times$(pos target $-$ pos) $+$ \code{damping}$\times$(vel
target $-$ vel). With \code{stiffness=0} only the second term is left, so
\srcpath{set_joint_velocity_target} is all the env needs to call
(Chapter~\ref{ch:env}). For the steering joint both gains are set and the
env calls \srcpath{set_joint_position_target}; its velocity target stays at the
default of zero, so the damping term acts as a brake on fast steering.
\end{isaacbox}

\section{The soil}

\codepart{\srcroot/luna_hifi_tasks/tricycle/tricycle_env_cfg.py}{114}{136}

\begin{explain}
\li{118} An \code{MPMObjectCfg} describes one particle body per env.
\li{119} Its prim under each env. Newton needs a prim to attach the particle
  cloud to.
\li{120} Default initial state: the particle box is not moved or rotated, so
  its local coordinates (lines 67--68) are env-local.
\li{121} The spawner: a regular lattice in a box.
\lis{122}{123} The box corners.
\li{124} The target voxel size; together with line 125 it fixes the lattice
  resolution.
\li{125} One particle per cell along each axis.
\li{126} Put particles at cell centres, so each particle's mass equals the
  cell volume times density and the total equals the box mass. The
  alternative, \code{"boundary"}, also places particles on the box faces.
\li{127} No random jitter: a perfectly regular lattice.
\li{128} Comment: the intent is firm ground, not loose sand.
\li{129} The material.
\li{130} Density 1800~kg/m$^3$, a dense regolith or packed soil. Each of the
  720 particles weighs $0.05^3 \times 1800 = 0.225$~kg; a strip is 162~kg.
\li{131} Friction coefficient $0.84 = \tan 40^\circ$: a steep internal
  friction angle. With zero \code{yield_stress} (cohesion) the material is
  a pure frictional granular medium.
\li{132} Yield pressure $10^{12}$~Pa: the soil never yields in compression,
  so wheels cannot compact it away. The defaults of the other fields
  (\code{young_modulus} $10^{15}$~Pa, \code{poisson_ratio} 0.3, no
  viscosity, damping, hardening or dilatancy) are inherited.
\li{133} End of material.
\li{134} Colour handed to the particle visualiser.
\lis{135}{136} Close spawner and cfg.
\end{explain}

\begin{warnbox}{parameter semantics changed from the old custom MPM}
In Newton's implicit MPM, \code{friction} is $\tan\phi$ and
\code{yield_stress} is cohesion in pascals; the yield surface is
$\tau_{\max}(p) = \texttt{yield\_stress} + \texttt{friction}\cdot(p-p_{\min})$.
Numbers from a Drucker--Prager $\alpha$ or a log-strain cohesion model do not
carry over; the runbook says ``recalibrate, don't port''. There are also no
\code{hardening_rate}/\code{softening_rate} fields.
\end{warnbox}

\section{The hidden slab}

\codepart{\srcroot/luna_hifi_tasks/tricycle/tricycle_env_cfg.py}{139}{161}

\begin{explain}
\li{139} A small factory function returning the slab config. A function
  rather than a constant keeps the long expression readable and mirrors the
  reference task.
\li{140} Docstring: the whole reason for the slab.
\li{141} A \code{RigidObjectCfg}: a single rigid body per env with no joints.
\li{142} Its prim, \code{MPMGround} under each env. This exact name is what
  the MPM coupler entry lists as its body (line 262).
\li{143} Placed at the slab centre from line 74.
\li{144} Spawn a box primitive.
\li{145} With the size from line 73.
\lis{146}{149} Rigid-body schema: enabled, and \emph{kinematic}. A kinematic
  body has a pose but no dynamics: it does not fall, and nothing can push it.
  It exists only to be a collider.
\lis{150}{154} Newton collision schema: collisions on; a
  \code{contact_margin} of half a voxel inflates the slab's collision
  surface outward by 2.5~cm so the MPM grid sees it early; \code{contact_gap}
  zero adds no extra detection distance.
\lis{155}{158} Friction of the slab against particles: static 0.8, dynamic
  0.7.
\li{159} Invisible: it is a physics-only helper and would otherwise hide the
  soil in the viewer.
\lis{160}{161} Close.
\end{explain}

\section{The scene}

\codepart{\srcroot/luna_hifi_tasks/tricycle/tricycle_env_cfg.py}{164}{183}

\begin{explain}
\li{169} A configclass \ldots
\li{170} \ldots\ inheriting \code{InteractiveSceneCfg}, which contributes
  \code{num_envs}, \code{env_spacing}, \code{replicate_physics} and friends.
  Every attribute added here becomes an entry of the scene.
\lis{171}{175} The ground plane. \code{AssetBaseCfg} is for things that
  have a prim but no per-env state; its path is absolute (one plane for all
  envs), spawned by \code{GroundPlaneCfg} (a 100$\times$100~m plane with a
  default material). \code{collision_group=-1} puts it in the global
  collision group so it collides with every env's bodies instead of only
  with one env. It belongs to the rigid coupler entry via
  \srcpath{include_static_shapes=True} (line 236); the soil never sees it.
\lis{176}{179} A dome light, so the OpenGL viewer is not dark. Colour and
  intensity are cosmetic.
\li{181} The car: the scene entry named \code{tricycle}, which is what
  \code{scene["tricycle"]} returns in the env.
\li{182} The slab, entry name \code{mpm_ground}.
\li{183} The soil, entry name \code{soil}.
\end{explain}

\begin{isaacbox}{scene entries and their order}
\code{InteractiveScene} iterates the cfg's fields in declaration order and
spawns each. Order barely matters here, but it decides the order in which
Newton is handed bodies, and the entry \emph{names} are the keys used by
\code{scene[...]} and by the reset (\srcpath{scene.reset(env_ids)} calls
\code{reset} on every asset).
\end{isaacbox}

\section{The environment config}

\codepart{\srcroot/luna_hifi_tasks/tricycle/tricycle_env_cfg.py}{186}{221}

\begin{explain}
\li{191} A configclass \ldots
\li{192} \ldots\ inheriting \code{DirectRLEnvCfg}, which declares the
  required fields (\code{decimation}, \code{episode_length_s}, \code{scene},
  \code{observation_space}, \code{action_space}) and defaults for the rest
  (\code{sim}, \code{seed}, noise models, \code{events}, \ldots).
\li{193} Two physics steps per policy step. With \code{dt = 1/100} the
  policy acts at 50~Hz.
\li{194} Episodes last 2~s. \code{DirectRLEnv} converts this to
  \code{max_episode_length} $= \lceil 2.0 / (0.01\times2) \rceil = 100$
  control steps.
\li{196} The simulation config: 10~ms physics step and a render every
  \code{decimation} physics steps, i.e.\ once per control step (rendering
  more often than the policy acts would be wasted work; rendering less often
  triggers a warning). The \code{physics} field is filled in later by
  \code{__post_init__}.
\lis{198}{202} The scene from above with its replication settings: 8 envs by
  default (the runbook trains with \code{--num_envs 4}), 2~m apart on the
  env grid, and \srcpath{replicate_physics=True} (spawn env 0 and clone it
  rather than spawning every env from scratch, which is much faster and is
  what Newton's per-world builder expects).
\li{204} Comment.
\li{205} Two actions. An \code{int} becomes \code{Box(-inf, inf, (2,))};
  the env clamps to $[-1,1]$ itself.
\li{206} Action scale for throttle: at $+1$ the rear wheels are asked to
  spin at 20~rad/s, about 2~m/s at 10~cm radius. (The runbook's
  ``known-good'' table lists 8.0 for this knob; the committed value is 20.0.
  Both have been run; lower values make the car slower and gentler on the
  coupling.)
\li{207} Action scale for steering: $\pm0.5$~rad, inside the joint's
  $\pm0.6$~rad limit.
\lis{209}{210} Comment listing the observation vector; the sum is 15.
\li{211} The observation size. Becomes \code{Box(-inf, inf, (15,))} under
  the \code{"policy"} key.
\li{212} No separate critic state; \code{0} means ``no state space''
  (Isaac Lab treats a falsy value as absent). The critic sees the same 15
  numbers.
\li{214} Comment.
\li{215} If the car drifts more than 0.5~m sideways it has left the 0.6~m
  wide strip (with a little slack because the wheels are 32~cm apart).
\li{216} Flip threshold on the $z$ component of gravity expressed in the
  car's body frame; see the orange box below.
\lis{218}{220} A named knob for the proxy mass scale so that it can be
  swept from the command line
  (\srcpath{env.proxy_mass_scale=5}) without editing the file. Line 276 reads
  it.
\end{explain}

\begin{warnbox}{projected gravity reads $+1$ upright for this asset}
\code{projected_gravity_b} is the world gravity direction rotated into the
body frame. For most Isaac Lab robots ``upright'' gives $(0,0,-1)$. For this
converted asset it gives $(0,0,+1)$, so the flip test in
\code{tricycle_env.py} is \code{projected_gravity_b[:, 2] < 0.3}: below
$0.3$ means the car has rolled more than about $73^\circ$ from upright. If
you ever swap the asset, print this quantity once before trusting the
termination.
\end{warnbox}

\section{\texttt{\_\_post\_init\_\_}: the physics}

\codepart{\srcroot/luna_hifi_tasks/tricycle/tricycle_env_cfg.py}{223}{286}

\begin{explain}
\li{223} \code{configclass} calls this after the instance is built. Setting
  the physics here rather than as a class attribute is deliberate: it can
  read \srcpath{self.proxy_mass_scale}, and it guarantees a fresh
  \code{NewtonCfg} per instance.
\li{224} Assign a Newton backend to \code{sim.physics}. This single
  assignment is what makes the whole run use Newton instead of PhysX.
\li{225} The solver is a \emph{coupler} holding two entries and one proxy
  mapping (Figure~\ref{fig:coupler}).
\li{226} The list of entries.
\lis{227}{228} Entry one, named \code{rover}.
\lis{229}{234} Its solver: MuJoCo-Warp with Newton-side contacts, the
  \code{implicitfast} integrator, and 128 constraint rows and 128 contacts
  per env (Chapter~\ref{ch:physics} explains each).
\li{235} The bodies this entry owns: every body under each env's
  \code{Tricycle} prim (descendants are included by the regex semantics),
  i.e.\ the whole car.
\li{236} Also own every shape that has no body: the ground plane. Without
  this the car would fall through the world floor.
\li{237} Two rigid sub-steps of 5~ms inside each 10~ms coupled step, for
  accuracy of the stiff wheel drives.
\li{238} End of entry one.
\lis{239}{240} Entry two, named \code{soil}.
\li{241} An implicit MPM solver.
\lis{242}{254} Grid and numerics, all explained in
  Section~\ref{sec:mpm}: 5~cm sparse grid without padding, P0 strain, APIC
  transfer, at most 24 iterations to $10^{-4}$, automatic warm start and
  solver choice, Q1 velocities, 27-point collider sampling with forward
  velocities, and one independent grid per env.
\lis{255}{256} Rule from the header: no particle push-out on a coupled
  entry.
\lis{257}{260} Sparse-grid capacities shared by all envs, in the required
  order active $\geq$ leaf $\geq$ lower $\geq$ upper:
  $16384 / 8192 / 4096 / 1024$. These are the 4-env, 6~GB numbers. For more
  envs raise all four roughly in proportion; \code{--num_envs} does not do it
  for you (Section~\ref{sec:step_cfg}).
\li{261} End of the MPM solver config.
\li{262} The only body the soil solver knows: the slab. Everything the soil
  should collide with must either be here or arrive as a proxy.
\li{263} This entry owns all particles in the model.
\li{264} It does not own the body-less shapes (the ground plane); those went
  to the rigid entry.
\li{265} It owns no joints (the slab has none).
\li{266} One MPM sub-step per coupled step.
\li{267} Step in place, mandatory for MPM.
\lis{268}{269} Close entry two and the list.
\li{270} The proxy mappings.
\li{271} One mapping.
\lis{272}{273} From the rigid entry into the soil entry.
\li{274} Which bodies to mirror: the three wheel bodies. The regex matches
  \code{rear_left_body}, \code{rear_right_body} and
  \srcpath{fork/front_wheel_body} because \code{.*} also matches the slash; it
  does not match \code{chassis} or \code{fork}, which never touch soil.
\li{275} Lagged mode: poses go over at the start of the step, impulses come
  back one step late (Section~\ref{sec:coupling}).
\li{276} The proxies weigh ten wheels each (0.5~kg $\times$ 10), from the
  knob on line 220.
\li{277} \code{None} means the soil entry reuses the shared outer contacts
  instead of running its own collision pipeline for the proxies.
\lis{278}{279} Close the mapping and the list.
\li{280} One relaxation pass per coupled step.
\li{281} End of the coupler.
\li{282} Newton's rigid collision pipeline, with the particle-vs-shape
  (``soft'') contact buffer disabled because the MPM solver handles particle
  contacts itself.
\li{283} One coupled solve per physics step.
\li{284} Capture the step as a CUDA graph.
\li{285} End of \code{NewtonCfg}.
\li{286} Let Newton execute the actuator drives natively.
\end{explain}

\section{\texttt{play\_mode}: settings for playback}

\codepart{\srcroot/luna_hifi_tasks/tricycle/tricycle_env_cfg.py}{288}{300}

\begin{explain}
\li{288} Isaac Lab calls this method on the env cfg when running
  \code{play} (Section~\ref{sec:play}). It never runs during training.
\li{289} The visualiser config is imported inside the method so that a
  headless training machine without the visualiser package can still import
  this file.
\li{291} The base implementation caps \code{num_envs} at 50 and switches
  off observation noise.
\li{292} Cap further to 4 envs: the soil grid capacities were sized for 4.
\li{293} Attach one visualiser to the simulation.
\li{294} Newton's OpenGL viewer.
\li{295} Draw the MPM particles. The default is \code{False}, which is why
  ``soil invisible'' is a runbook troubleshooting entry.
\li{296} Draw them in the soil colour.
\lis{297}{298} Camera position and target, in metres, looking at the middle
  of the strip from front-left and above.
\lis{299}{300} Close.
\end{explain}
